\documentclass[12pt, titlepage]{article}

\usepackage{booktabs}
\usepackage{tabularx}
\usepackage{hyperref}
\hypersetup{
    colorlinks,
    citecolor=blue,
    filecolor=black,
    linkcolor=red,
    urlcolor=blue
}
\usepackage[round]{natbib}

\input{../Comments.text}
\input{../Common.text}

\begin{document}

\title{System Verification and Validation Plan for \progname{}} 
\author{\authname}
\date{\today}
	
\maketitle

\pagenumbering{roman}

\section*{Revision History}

\begin{tabularx}{\textwidth}{p{3cm}p{2cm}X}
\toprule {\bf Date} & {\bf Version} & {\bf Notes}\\
\midrule
Date 1 & 1.0 & Notes\\
Date 2 & 1.1 & Notes\\
\bottomrule
\end{tabularx}

~\\
\wss{The intention of the VnV plan is to increase confidence in the software.
However, this does not mean listing every verification and validation technique
that has ever been devised.  The VnV plan should also be a \textbf{feasible}
plan. Execution of the plan should be possible with the time and team available.
If the full plan cannot be completed during the time available, it can either be
modified to ``fake it'', or a better solution is to add a section describing
what work has been completed and what work is still planned for the future.}

\wss{The VnV plan is typically started after the requirements stage, but before
the design stage.  This means that the sections related to unit testing cannot
initially be completed.  The sections will be filled in after the design stage
is complete.  the final version of the VnV plan should have all sections filled
in.}

\newpage

\tableofcontents

\listoftables
\wss{Remove this section if it isn't needed}

\listoffigures
\wss{Remove this section if it isn't needed}

\newpage

\section{Symbols, Abbreviations, and Acronyms}

\renewcommand{\arraystretch}{1.2}
\begin{tabular}{l l} 
  \toprule		
  \textbf{symbol} & \textbf{description}\\
  \midrule 
  T & Test\\
  \bottomrule
\end{tabular}\\

\wss{symbols, abbreviations, or acronyms --- you can simply reference the SRS
  \citep{SRS} tables, if appropriate}

\wss{Remove this section if it isn't needed}

\newpage

\pagenumbering{arabic}

This document ... \wss{provide an introductory blurb and roadmap of the
  Verification and Validation plan}

\section{General Information}

\subsection{Summary}

PronouncedFish is an assistive tool intended for SLPs to streamline their workflow in diagnosing and treating patients with speech sound disorders (SSDs). It can help manage treatment sessions and offer complementary insights to assist with their consultations. The important features that PronouncedFish provides are real-time analysis of patient speech to detect SSDs at the phonemic level; Management, storage, and analysis of treatment sessions; providing practice material for patients to use in a session; and tracking of patient speech data overtime that can be used to guide the focus of future treatment sessions.

\subsection{Objectives}

With this VnV Plan, Ghoti aims to discuss how it will validate the correctness, quality, robustness, and efficiency of PronouncedFish. With this comes the objectives that Ghoti wants to achieve for this document. 

\subsubsection{Full Functionality of the System}
Each component of the system we implement in code needs to be tested for correctness. These components include:
\begin{itemize}
    \item Audio Management: The audio management component should be tested to see if it can take the audio provided by a patient and process it to be as clear as possible to be used by other components in the system. (FR-AC1, FR-AC2, FR-AC3)
    \item Audio Analysis: The audio analysis component should be tested to see if it can take a piece of audio, usually some kind of sentence, and return a set of predicted phonemes that corresponds to the productions made by the patient. Should the analysis fail, notifications of failure should be broadcast throughout the system and to the SLP. (FR-SA1, FR-SA2, FR-SA3, FR-AC3)
    \item Session Management: The session management component should be tested to see if it can create, delete, and modify sessions. It should also be able to take data from a session and provide it to the rest of the system for later analysis by the SLP. (FR-SM1, FR-SM2, FR-SM3, FR-SM4, FR-DMR6, FR-DMR7)
    \item Practice Item Management: The practice item management component should be tested to see if it can create, delete, and modify practice materials that a patient would use in a session. It should also only respond to SLPs authorized with the system and software maintenance teams. (FR-PAR1, FR-PAR2, FR-PAR3, FR-PAR4, FR-PAR5, FR-PAR6, FR-PAR7)
    \item Progress Tracking: The progress tracking component should be tested to see if it can store patient data to later be used by SLPs for analysis. (FR-DMR1, FR-DMR6, FR-DMR8, FR-DMR9, FR-DMR10, FR-DMR11)
    \item Data Visualization: The data visualization component should be tested to see if it can take existing data for a patient and output it in a visual, human-readable format for SLPs. (FR-RDV1, FR-RDV2, FR-RDV3, FR-RDV4, FR-RDV5, FR-RDV6, FR-RDV7, FR-RDV8, FR-RDV9)
\end{itemize}

\subsubsection{Personal Health Information Protection Act (PHIPA) Compliance}
The VnV plan should ensure that each component is within compliance of PHIPA to protect sensitive personal health information (PHI). PHIPA compliance comes with a few categories:
\begin{itemize}
    \item Data Encryption: The system should ensure that any PHI that is stored or passed around in the system must be encrypted to prevent sensitive data from being exposed. (NFR-SAP1, NFR-SAP2)
    \item The system should ensure that any time PHI is accessed, edited, or deleted, an immutable log of when the data was accessed is stored in the system for traceability. (FR-DMR4, FR-DMR5)
    \item The system should ensure that anywhere PHI is stored, it can only be accessed by authorized users like an SLP. (NFR-USA4, NFR-LAL1)
\end{itemize}

\subsubsection{High Accuracy of the Phonemic Identification Model}

For the model we use for phonemic identification within the system, experts in the field including our supervisors have mentioned that the model should be 80\% accurate. (NFR-ACCU1)

\subsubsection{High Responsiveness of the System}

The system should have a response time of less than or equal to 10 seconds. 10 seconds was chosen as our metric due to this paper by \cite{software_response_times} which mentions that most people will on average have an attention span of about 10 seconds before wanting to perform other tasks. Multiple elements of the system will have a strong impact on the response time including the processing time of the phonemic identification model and communications with local servers. The tests will ensure that any action taken by the user will fall within this 10 second response time. (NFR-PER1, NFR-REL1)

\subsubsection{Great User Experience for Patients}

User experience is a key focus in our project since some of the primary stakeholders of PronouncedFish are children. If the patient gets bored, frustrated, or upset by the system it can have a significant impact on the effectiveness of the session. Thus, we want user ratings of PronouncedFish to be high and frustrations low (NFR-USA1, NFR-USA3, NFR-USA4, FR-PAR7, FR-RDV8)

\subsubsection{Portability Across Hardware and Operating Systems}

Stakeholders will not all be using the same hardware, therefore, there is a need for the system to work across multiple kinds of hardware and operating systems. Our tests and software should work with most popular operating systems (Linux, Windows, iOS, macOS, android) and on multiple different kinds of hardware (laptop, desktop, iPad, android device, etc.). This will be tested manually through running PronouncedFish on a variety of devices. (NFR-POR1, NFR-POR2, NFR-POR3)

\subsection{Out-of-Scope Objectives}

\subsubsection{Long-Term Maintenance Testing}

The system will require updates and long-term maintenance to run at optimal capacity. Unfortunately for the duration of this capstone, Ghoti does not have an extensive period of time to properly test how regular maintenance and updates would look in a real setting. Assumed testing would be update and maintenance logs showing what parts of the system were updated and at what time. (NFR-MAI3, NFR-REL4)

\subsubsection{External Library Functionality Testing}

The system will be using external libraries to function; however, we assume the libraries are always providing the services that they advertised and do not need to be independently checked for successful functionality. Assumed testing would be to go through every possible unit path in the provided services from the library.

\subsection{Custom Documents}

\subsubsection{User/Stakeholder Interview Report}

The User/Stakeholder Interview Report serves as an important document for visualizing SLP workflows and eliciting specific requirements for PronouncedFish. By conducting structured interviews with our supervisor SLPs, we ensure that the system’s design follows clinical best practices and integrates smoothly with an SLP's daily operations. These key insights will help us ensure core functionalities are met and follow the strict rules from PHIPA surrounding the protection of PHI. (FR-SM1, FR-SA1, FR-PAR1, FR-DMR8, NFR-LAL1)

\subsubsection{Usability Report}

The Usability Report is designed to validate the system's user experience (UX) before the completion of the capstone. After developing a functional prototype, the students of Ghoti's supervisors will be able to experiment with using PronouncedFish. The primary goal of this document is not just feedback collection, but quantifying key performance indicators such as response time and the accuracy of the machine learning model when identifying phonemes. This ensures that the design meets high standards of intuitive functionality, validating our objective to provide a positive user experience critical for clinical adoption. (NFR-USA1, NFR-USA3, NFR-PER1, NFR-ACCU1)

\subsection{Relevant Documentation}

This section outlines the primary documents that govern the system's scope, constraints, and core functional requirements. These documents act as foundational references against which all development progress and testing validation are measured, ensuring comprehensive traceability between goals, requirements, and implementation. 

\subsubsection{Software Requirements Specification (SRS)}

The SRS serves as the formal blueprint for the entire application and is the definitive source of truth for all functional and non-functional system expectations. All requirements being tested will be sourced from here.

\subsubsection{Development Plan}

Details the project roadmap, risk assessment, technology stack selection (libraries and frameworks), and overall architectural approach used in the system's implementation. (This document establishes the context for performance and feasibility testing.)

\subsubsection{Module Interface Specification (MIS)}

Provides fine-grained technical design details for core modules, defining specific access functions, execution sequences, and serves as primary source of information for creating unit tests.

\section{Plan}

\wss{Introduce this section.  You can provide a roadmap of the sections to
  come.}

\subsection{Verification and Validation Team}

\wss{Your teammates.  Maybe your supervisor.
  You should do more than list names.  You should say what each person's role is
  for the project's verification.  A table is a good way to summarize this information.}

\subsection{SRS Verification}

\wss{List any approaches you intend to use for SRS verification.  This may
  include ad hoc feedback from reviewers, like your classmates (like your
  primary reviewer), or you may plan for something more rigorous/systematic.}

\wss{If you have a supervisor for the project, you shouldn't just say they will
read over the SRS.  You should explain your structured approach to the review.
Will you have a meeting?  What will you present?  What questions will you ask?
Will you give them instructions for a task-based inspection?  Will you use your
issue tracker?}

\wss{Maybe create an SRS checklist?}

\subsection{Design Verification}

\wss{Plans for design verification}

\wss{The review will include reviews by your classmates}

\wss{Create a checklists?}

\subsection{Verification and Validation Plan Verification}

\wss{The verification and validation plan is an artifact that should also be
verified.  Techniques for this include review and mutation testing.}

\wss{The review will include reviews by your classmates}

\wss{Create a checklists?}

\subsection{Implementation Verification}

\wss{You should at least point to the tests listed in this document and the unit
  testing plan.}

\wss{In this section you would also give any details of any plans for static
  verification of the implementation.  Potential techniques include code
  walkthroughs, code inspection, static analyzers, etc.}

\wss{The final class presentation in CAS 741 could be used as a code
walkthrough.  There is also a possibility of using the final presentation (in
CAS741) for a partial usability survey.}

\subsection{Automated Testing and Verification Tools}

\wss{What tools are you using for automated testing.  Likely a unit testing
  framework and maybe a profiling tool, like ValGrind.  Other possible tools
  include a static analyzer, make, continuous integration tools, test coverage
  tools, etc.  Explain your plans for summarizing code coverage metrics.
  Linters are another important class of tools.  For the programming language
  you select, you should look at the available linters.  There may also be tools
  that verify that coding standards have been respected, like flake9 for
  Python.}

\wss{If you have already done this in the development plan, you can point to
that document.}

\wss{The details of this section will likely evolve as you get closer to the
  implementation.}

\subsection{Software Validation}

\wss{If there is any external data that can be used for validation, you should
  point to it here.  If there are no plans for validation, you should state that
  here.}

\wss{You might want to use review sessions with the stakeholder to check that
the requirements document captures the right requirements.  Maybe task based
inspection?}

\wss{For those capstone teams with an external supervisor, the Rev 0 demo should 
be used as an opportunity to validate the requirements.  You should plan on 
demonstrating your project to your supervisor shortly after the scheduled Rev 0 demo.  
The feedback from your supervisor will be very useful for improving your project.}

\wss{For teams without an external supervisor, user testing can serve the same purpose 
as a Rev 0 demo for the supervisor.}

\wss{This section might reference back to the SRS verification section.}

\section{System Tests}

\wss{There should be text between all headings, even if it is just a roadmap of
the contents of the subsections.}

\subsection{Tests for Functional Requirements}

\wss{Subsets of the tests may be in related, so this section is divided into
  different areas.  If there are no identifiable subsets for the tests, this
  level of document structure can be removed.}

\wss{Include a blurb here to explain why the subsections below
  cover the requirements.  References to the SRS would be good here.}

\subsubsection{Area of Testing1}

\wss{It would be nice to have a blurb here to explain why the subsections below
  cover the requirements.  References to the SRS would be good here.  If a section
  covers tests for input constraints, you should reference the data constraints
  table in the SRS.}
		
\paragraph{Title for Test}

\begin{enumerate}

\item{test-id1\\}

Control: Manual versus Automatic
					
Initial State: 
					
Input: 
					
Output: \wss{The expected result for the given inputs.  Output is not how you
are going to return the results of the test.  The output is the expected
result.}

Test Case Derivation: \wss{Justify the expected value given in the Output field}
					
How test will be performed: 
					
\item{test-id2\\}

Control: Manual versus Automatic
					
Initial State: 
					
Input: 
					
Output: \wss{The expected result for the given inputs}

Test Case Derivation: \wss{Justify the expected value given in the Output field}

How test will be performed: 

\end{enumerate}

\subsubsection{Area of Testing2}

...

\subsection{Tests for Nonfunctional Requirements}

\subsubsection{Accuracy and responsiveness}

Use a consistent evaluation set so measurements can be compared across
model changes, as required by the Development Plan.

\paragraph{Measure phoneme detection accuracy}
\begin{enumerate}
    \item \textbf{ST-N01}\\

    Type: Dynamic, Automatic

    Initial State: A set of labelled children's speech clips from TalkBank
    (e.g.\ PhonBank), with reference phoneme transcriptions, is loaded.
    Clips used for tuning are excluded.

    Input/Condition: Run every clip through the recognizer and compare the
    output phonemes against the labelled reference phonemes.

    Output/Result: Phoneme identification accuracy is at least 87\%
    (NFR-ACCU1; FR-SA1, FR-SA2, FR-SA3). Report accuracy overall and per
    target phoneme.

    How test will be performed: Align recognized and reference phonemes for
    each clip. Accuracy is the number of correctly identified phonemes divided
    by the total number of reference phonemes. Rerun the same clip set after
    each model change.
\end{enumerate}

\paragraph{Measure processing time from audio capture to phoneme estimation}
\begin{enumerate}
    \item \textbf{ST-N02}\\

    Type: Dynamic, Automatic measurement

    Initial State: The device and software are
    configured to record speech and process it.

    Input/Condition: Input audio file for a particular word.

    Output/Result: For each case, phonemic estimation results are
    visible in less than 10 seconds (NFR-PER1). Record duration for each processing completion and categorize by the specific phrase. Report all observations plus processing time divided by audio
    duration, median, and maximum.

    How test will be performed: Measure from completed audio
    capture/submission to rendered result, including transfer and rendering,
    with at least 30 runs per representative condition as a planned sampling
    choice. Log hardware, network, clip duration and model size. Resolve
    whether capture time belongs in the requirement's timing boundary before
    acceptance. Basis: NFR-PER1 sets a strict less than 10 seconds criterion;
    Development Plan Section~10.5 adds real-time factor and end to end
    latency.
\end{enumerate}

\subsubsection{Usability and patient protection}

Use SLP task observations for navigation and efficiency. Assess the child-facing boundary
separately from clinician access.

\paragraph{Learn the workflow within one hour}
\begin{enumerate}
    \item \textbf{ST-N03}\\

    Type: Dynamic, Manual

    Initial State: An SLP unfamiliar with the application has standard
    onboarding material and synthetic data.

    Input/Condition: Within a one hour session, ask the SLP to configure
    practice, record audio and analyze it, override a result, inspect
    charts and observe patient history.

    Output/Result: The SLP completes the core navigation tasks independently
    within one hour of use (NFR-USA1). Record task times, assistance, failures
    and comments.

    How test will be performed: The UI/UX Lead observes without guiding the
    final task run. They record mistakes, time to complete each task, and any questions or comments.
    The SLP is asked to provide feedback on the experience.
\end{enumerate}

\begin{enumerate}
    \item \textbf{ST-N04}\\

    Type: Dynamic usability and recovery checks

    Initial State: A test SLP account exists with synthetic critical data
    (for example, patient records and session recordings).

    Input/Condition: As the SLP, attempt to delete each type of critical data
    item. Where a confirmation prompt appears, test both cancelling and
    confirming. Where deletion is allowed without a prompt, attempt recovery
    at several points within 72 hours and again after 72 hours have elapsed.

    Output/Result: No critical data item is permanently deleted without
    either a confirmation step or a recovery option (NFR-USA2). Cancelling a
    confirmation leaves the data unchanged. Deleted data is fully restored
    when recovery is requested within 72 hours, with contents identical to
    the original.

    How test will be performed: The Development Lead and an independent
    reviewer walk through each data type, automated script compares restored contents
    against the original synthetic data, and check the system clock or
    retention settings to verify the 72-hour window. Record any data type
    that can be lost without a safeguard; review is blocked if the critical
    data list is absent. Basis: NFR-USA2.
\end{enumerate}

\paragraph{Compare workflow time with manual tools}
\begin{enumerate}
    \item \textbf{ST-N05}\\
    Type: Dynamic, Manual comparative measurement

    Initial State: An SLP has learned the application and has equivalent paper
    or spreadsheet tools.

    Input/Condition: Complete matched session and after-session documentation
    tasks with identical attempts using software and manual tools.

    Output/Result: Total software task time is lower than manual task time,
    with equivalent completed records (NFR-USA3). Report paired times,
    corrections and assistance; a faster incomplete record does not pass.

    How test will be performed: The UI/UX Lead checks
    where multiple participants are available and records session and
    after-session times separately and together. Report sample size and
    learning effects. Basis: NFR-USA3 requires less time than equivalent manual
    tools, including work outside the session.
\end{enumerate}

\paragraph{Restrict the child-facing experience}
\begin{enumerate}
    \item \textbf{ST-N06}\\

    Type: Dynamic, Mixed manual and automatic

    Initial State: Child practice mode is active; clinician functions require
    the configured access boundary.

    Input/Condition: Try visible controls, back navigation, direct clinician
    URLs and direct API requests from the child context.

    Output/Result: The child can use the gamified experience but cannot access
    patient records, analytics administration or clinician controls beyond
    that experience (NFR-USA4). Direct navigation and API calls do not bypass
    restrictions.

    How test will be performed: The Systems Architect/Development Lead tests
    each route with the UI/UX Lead and captures denied requests and
    screenshots. Verify returning to clinician mode requires the configured
    authorization. Basis: NFR-USA4 limits the child's access to the practice
    experience.
\end{enumerate}

\subsubsection{Maintainability and portability}

Checks supported hardware and software environments. Check deployment, patching, and extension.

\begin{enumerate}
    \item \textbf{ST-N07}\\

    Type: Dynamic performance check

    Initial State: A clean build environment and a prepared patch exist, and
    the approved build and deployment pipeline is available.

    Input/Condition: Apply a representative patch and run the pipeline from
    build start to a completed deployment, timing each stage.

    Output/Result: Building and deploying the patch completes successfully in
    less than 30 minutes (NFR-MAI3).

    How test will be performed: The Development Lead and an independent
    reviewer time the pipeline run from pipeline logs and repeat it three
    times. Record durations and any failures; review is blocked if the
    pipeline is unavailable. Basis: NFR-MAI3.
\end{enumerate}

\paragraph{Verify platform and browser compatibility}
\begin{enumerate}
    \item \textbf{ST-N08}\\

    Type: Dynamic, Manual

    Initial State: A versioned compatibility matrix is recorded with physical
    devices for microphone checks.

    Input/Condition: Exercise Windows with Chrome, Firefox and Edge; macOS with
    Safari, Chrome and Firefox; iPad/iOS with Safari; and Android tablet with
    Chrome. Include desktop, laptop and tablet form factors.

    Output/Result: Setup, microphone permission and capture, analysis, SLP
    controls, IPA, charts and history operate without blocking layout or
    functional defects in every tested combination (NFR-POR1, NFR-POR2,
    NFR-POR3; FR-RDV9).

    How test will be performed: Repeat a compact ST-F01/F03/F08/F09/F11
    workflow in each matrix entry; inspect orientation and touch controls on
    tablets; test microphone functionality. Add any other claimed supported combination to the matrix before
    declaring support. Basis: The SRS lists four operating systems, four
    browsers and three device types; the matrix is a representative test plan,
    not proof of every OS-browser combination.
\end{enumerate}

\subsubsection{Reliability, security and release review}

These checks address operational hazards, including data loss, and unauthorized access.

\paragraph{Verify operating periods and maintenance scheduling}
\begin{enumerate}
    \item \textbf{ST-N09}\\

    Type: Mixed Dynamic and Static

    Initial State: Deployment operating timezone, full and minimal service
    definitions, load and maintenance policy are approved.

    Input/Condition: Exercise login, storage, retrieval and analysis around
    06:00 and 00:00; run an accelerated schedule check and inspect the
    maintenance calendar.

    Output/Result: The approved full service level operates from 06:00 to
    00:00 and the approved minimum remains available from 00:00 to 06:00.
    Routine maintenance is scheduled in minimal hours at the approved
    interpretation of the one to two month interval (NFR-REL2, NFR-REL3,
    NFR-REL4).
    Check server latency and ensure it is under 10 seconds for this period (NFR-REL1).

    How test will be performed: Use a controlled clock and scheduler
    inspection, then planned observation under the approved workload.
\end{enumerate}

\paragraph{Protect data and secure communications}
\begin{enumerate}
    \item \textbf{ST-N10}\\

    Type: Static inspection and Dynamic security checks

    Initial State: The approved PHI inventory and encryption baseline identify
    every component and storage location.

    Input/Condition: Create and transfer synthetic PHI; inspect browser-to-API,
    service-to-service and database traffic, persistent stores, logs,
    temporary files and backups.

    Output/Result: PHI is encrypted from collection through its approved
    storage lifecycle; intercomponent traffic uses the approved secure
    protocol (NFR-SAP1, NFR-SAP2; NFR-LAL1). No unintended plaintext PHI
    appears in inspected storage or traffic. Approved key access and
    certificate validation controls are active.

    How test will be performed: The Systems Architect/Development Lead and an
    independent team reviewer inspect configuration, packet traces and
    synthetic storage samples. Test denial of an untrusted certificate and an
    insecure endpoint. Record coverage and gaps; review is blocked if the
    baseline is absent. Basis: NFR-SAP1 and NFR-SAP2 require encryption; exact
    cryptographic acceptance comes from an approved security baseline, not this
    plan.
\end{enumerate}

\subsection{Legal and Licensing}

Checks to ensure that the system is PHIPA compliant and open source.

\begin{enumerate}
    \item \textbf{ST-N11}\\

    Type: Static inspection and Dynamic security checks

    Initial State: The approved PHIPA checklist and deployment region list
    exist.

    Input/Condition: Process synthetic patient records through storage,
    transfer and access paths; attempt access with unauthorized roles; review
    the repository license and dependency manifest.

    Output/Result: Synthetic PHI is stored, transferred and accessed only as
    the PHIPA checklist permits, and stays in approved residency locations
    (NFR-LAL1). Unauthorized access is denied. The repository includes the
    GNU GPL text and all dependencies are GPL-compatible (NFR-LAL2).

    How test will be performed: The Development Lead and an independent
    reviewer compare configuration, storage locations and access logs against
    the checklist, and run a dependency license scan. Record findings; review
    is blocked if the checklist is absent. Basis: NFR-LAL1 and NFR-LAL2; legal
    sign-off is outside this plan.
\end{enumerate}



\subsubsection{Area of Testing1}
		
\paragraph{Title for Test}

\begin{enumerate}

\item{test-id1\\}

Type: Functional, Dynamic, Manual, Static etc.
					
Initial State: 
					
Input/Condition: 
					
Output/Result: 
					
How test will be performed: 
					
\item{test-id2\\}

Type: Functional, Dynamic, Manual, Static etc.
					
Initial State: 
					
Input: 
					
Output: 
					
How test will be performed: 

\end{enumerate}

\subsubsection{Area of Testing2}

...

\subsection{Traceability Between Test Cases and Requirements}

\wss{Provide a table that shows which test cases are supporting which
  requirements.}

\section{Unit Test Description}

\wss{This section should not be filled in until after the MIS (detailed design
  document) has been completed.}

\wss{Reference your MIS (detailed design document) and explain your overall
philosophy for test case selection.}  

\wss{To save space and time, it may be an option to provide less detail in this section.  
For the unit tests you can potentially layout your testing strategy here.  That is, you 
can explain how tests will be selected for each module.  For instance, your test building 
approach could be test cases for each access program, including one test for normal behaviour 
and as many tests as needed for edge cases.  Rather than create the details of the input 
and output here, you could point to the unit testing code.  For this to work, you code 
needs to be well-documented, with meaningful names for all of the tests.}

\subsection{Unit Testing Scope}

\wss{What modules are outside of the scope.  If there are modules that are
  developed by someone else, then you would say here if you aren't planning on
  verifying them.  There may also be modules that are part of your software, but
  have a lower priority for verification than others.  If this is the case,
  explain your rationale for the ranking of module importance.}

\subsection{Tests for Functional Requirements}

\wss{Most of the verification will be through automated unit testing.  If
  appropriate specific modules can be verified by a non-testing based
  technique.  That can also be documented in this section.}

\subsubsection{Module 1}

\wss{Include a blurb here to explain why the subsections below cover the module.
  References to the MIS would be good.  You will want tests from a black box
  perspective and from a white box perspective.  Explain to the reader how the
  tests were selected.}

\begin{enumerate}

\item{test-id1\\}

Type: \wss{Functional, Dynamic, Manual, Automatic, Static etc. Most will
  be automatic}
					
Initial State: 
					
Input: 
					
Output: \wss{The expected result for the given inputs}

Test Case Derivation: \wss{Justify the expected value given in the Output field}

How test will be performed: 
					
\item{test-id2\\}

Type: \wss{Functional, Dynamic, Manual, Automatic, Static etc. Most will
  be automatic}
					
Initial State: 
					
Input: 
					
Output: \wss{The expected result for the given inputs}

Test Case Derivation: \wss{Justify the expected value given in the Output field}

How test will be performed: 

\item{...\\}
    
\end{enumerate}

\subsubsection{Module 2}

...

\subsection{Tests for Nonfunctional Requirements}

\wss{If there is a module that needs to be independently assessed for
  performance, those test cases can go here.  In some projects, planning for
  nonfunctional tests of units will not be that relevant.}

\wss{These tests may involve collecting performance data from previously
  mentioned functional tests.}

\subsubsection{Module ?}
		
\begin{enumerate}

\item{test-id1\\}

Type: \wss{Functional, Dynamic, Manual, Automatic, Static etc. Most will
  be automatic}
					
Initial State: 
					
Input/Condition: 
					
Output/Result: 
					
How test will be performed: 
					
\item{test-id2\\}

Type: Functional, Dynamic, Manual, Static etc.
					
Initial State: 
					
Input: 
					
Output: 
					
How test will be performed: 

\end{enumerate}

\subsubsection{Module ?}

...

\subsection{Traceability Between Test Cases and Modules}

\wss{Provide evidence that all modules have been considered.}
				
\bibliographystyle{plainnat}

\bibliography{../../refs/References}

\newpage

\section{Appendix}

This is where you can place additional information.

\subsection{Symbolic Parameters}

The definition of the test cases will call for SYMBOLIC\_CONSTANTS.
Their values are defined in this section for easy maintenance.

\subsection{Usability Survey Questions?}

\wss{This is a section that would be appropriate for some projects.}

\newpage{}

\section{Appendix --- Generative AI Use Disclosure}

\wss{ChatGPT (version 5) was used to brainstorm ideas for the template for this disclose section.}

\subsection{AI Tools Used}

\wss{Give the name of the AI tool(s) used and the version.}

\subsection{How and Where Used}

\wss{Explain what the tool(s) were used for.  Examples include: brainstorming,
explaining a technical concept, reviewing the document, generating or modifying
text, generating or modifying diagrams, identifying errors or omissions or
inconsistencies.}

\wss{For each of the uses, identify which parts of the document were affected.}

\wss{You do not need to report every interaction, but you should disclose the 
uses that materially contributed to your work}

\subsection{Verification of AI-Generated Content}

\wss{\begin{itemize}
    \item Describe how the team checked AI-generated or AI-assisted material
    before including it in the document.
    \item In particular, verify factual claims, technical assertions,
    requirements, calculations, references, and other information for which an
    AI system may produce incorrect or fabricated results.
    \item \textbf{The team is responsible for the accuracy, completeness,
consistency, and appropriateness of the submitted document, regardless of
whether material was generated by a human or an AI system.}
\end{itemize}}

\newpage{}

\section{Appendix --- Reflection}

\wss{This section is not required for CAS 741}

The information in this section will be used to evaluate the team members on the
graduate attribute of Lifelong Learning.

\input{../Reflection.text}

\begin{enumerate}
  \item What went well while writing this deliverable? 
  \item What pain points did you experience during this deliverable, and how
    did you resolve them?
  \item What knowledge and skills will the team collectively need to acquire to
  successfully complete the verification and validation of your project?
  Examples of possible knowledge and skills include dynamic testing knowledge,
  static testing knowledge, specific tool usage, Valgrind etc.  You should look to
  identify at least one item for each team member.
  \item For each of the knowledge areas and skills identified in the previous
  question, what are at least two approaches to acquiring the knowledge or
  mastering the skill?  Of the identified approaches, which will each team
  member pursue, and why did they make this choice?
\end{enumerate}

\end{document}